% !PRINTED none
% Part-title / section-opening page: no running head / page number visible
% in any of the three bands (top band is essentially blank). Content:
% the "Achter Abschnitt." part title with its descriptive subtitle, a
% dividing rule, then the "§ 78." section heading and the opening of its
% body text.
\begin{center}
\gesperrt{Eighth Chapter.}

\gesperrt{The Group of the Transformation Equations and the Equation of Degree 5.}
\end{center}

\noindent\rule{0.3\textwidth}{0.4pt}

\begin{center}
§ 78.

\gesperrt{The Galois Group of the Transformation Equations for a Prime Degree.}
\end{center}

A more thorough algebraic study of the transformation equations requires a
knowledge of their \gesperrt{Galois} group. Since we have derived the transformation
equations from the division equations, whose group is known to us (§ 63), we
can form the group of the transformation equations as well.

The transformation equations arose (§ 65) from the fact that the roots of the
division equations admitted of a partition into rows which are not torn apart
by the permutations of the group of the division equation, but are only
interchanged among one another.

To each of these rows there corresponds a definite root of a transformation
equation, and the \gesperrt{group of the latter
therefore consists of the totality of the permutations that are produced among the rows by the group of the division equation.}

If the degree of transformation $n$ is \gesperrt{an odd prime $p$}, this group admits a
very elegant representation which is suited to further investigations, and we
now hold fast to this assumption.

In § 65, (3) the necessary and sufficient condition was already determined for
two roots of the division equation

\weq{1}{x_{\mu,\mu'},\quad x_{\nu,\nu'}}

to belong to the same row $R$, namely the congruence

\weq{2}{\mu\,\nu'-\nu\,\mu' \equiv 0 \pmod{p}.}
